\documentclass{article}[10pt]
\usepackage[utf8]{inputenc}
\usepackage{amssymb, amsmath, graphicx}
\usepackage{graphicx}
\usepackage{epsfig}
\usepackage{fancyhdr}
\usepackage{enumitem,mathtools}
\usepackage{a4wide}
\usepackage{amsthm}
\usepackage{geometry}
\usepackage{tikz,tikz-cd}
\usepackage[backgroundcolor=gray!8]{mdframed}

\newtheorem{theorem}{Theorem}

\theoremstyle{definition}
\newtheorem{exercise}[theorem]{Exercise}
\newtheorem*{solution*}{Solution}


%\newcommand{\mc}[1]{\mathcal{#1}}
%\newcommand{\NN}{\mathbb{N}}
%\newcommand{\p}[1]{\left( #1 \right)}
%\def\ol{\overline}
\newcommand{\bbH}{\mathbb{H}}
\newcommand{\bbR}{\mathbb{R}}
\newcommand{\Z}{\mathbb{Z}}


\input{../../exrepo2/definitions}

\input{../../exrepo2/mcm}
 \hidesolution


\begin{document}

\intro{%
    \noindent
    \rlap{Prof. Dr. A. Iozzi\\~\\}\hfill Symmetric Spaces \hfill\llap{FS 2026} \\~\\~\\
    {\centerline{\LARGE \textbf{Solutions Exercise Sheet 1}}}
    ~\\

  \hrule\bigskip%
}%

%%%%%%%%%%%% QUESTION 1 %%%%%%%%%%%%%%%
	
	\begin{enumerate}\item[1.] Let $G=\SL_{2}(\mathbb{R})$. The aim of this exercise is to show that $G_{\mathbb{Z}}=\SL_{2}(\mathbb{Z})$ is a lattice in $G$.
	\begin{enumerate}
		\item Argue that $G_{\mathbb{Z}}$ is discrete in $G$ and that both $G$ and $G_{\mathbb{Z}}$ are unimodular.
	\end{enumerate}
	
	%%%%%%% 1a solution
	
	\explanation{\begin{mdframed}
	\begin{solution*}
	$\SL(n,\Z)$ is contained in $\Z^{n\times n}$, which is a discrete subset of $\bbR^{n\times n}$; it follows that $G_\Z$ is discrete in $G$ by restriction from $\bbR^{n\times n}$ to $G$. In addition, $G_\Z$ is unimodular since it is discrete (and so the Haar measure is the counting measure).
	 For $G$ we recall the fact that $|\det(x_{ij})|^{-1}\,dx_{11}\cdots dx_{nn} $ is the bi-invariant Haar measure in $\operatorname{GL}(n,\bbR)$ and that $G$ is a closed normal subgroup of $\operatorname{GL}(n,\bbR)$ (see the Lie Groups lecture notes, Proposition 2.3, on page 28).
	\end{solution*}
	\end{mdframed}
}
From this we know that $G/G_{\mathbb{Z}}$ admits a nonzero $G$-invariant measure $\mu$ which is unique up to a non-zero constant. In order to show that $G_{\mathbb{Z}}$ is a lattice we have to show that $\mu(G/G_{\mathbb{Z}})<\infty$. For this, we will use the following fact:
		
\begin{enumerate}[label=(\alph*),start=2]
		\item Assume that there exists a measurable set $A\subseteq G$ of finite measure such that every $G_{\mathbb{Z}}$-orbit intersects $A$ (that is, for every $g\in G$ there exists some $\gamma\in G_{\mathbb{Z}}$ such that $g\gamma\in A$). Show that $\mu(G/G_{\mathbb{Z}})$ is finite.
	%%%%%%% 1b solution
	
	\explanation{\begin{mdframed}
	\begin{solution*}
		Recall Weil's formula (Lie Groups notes, Theorem 2.4 on page 34) which for the (integrable) characteristic function $\chi_A$ of $A$ in $G$ states that
		\begin{equation*}
			\mu(A) = \int_G\chi_A(g)\, dg = \int_{G/G_\Z}\left(\int_{G_\Z}\chi_A(\gamma h)\, dh\right) \, d(\gamma G_\Z)
		\end{equation*}
		By the assumption the inner integral is always greater than some absolute constant depending only on the Haar measure of $G_\Z$, and so we infer that $\mu(A) \geq c\mu(G/G_{\Z})$.
	\end{solution*}
	\end{mdframed}}
	
		\item  Show that the map sending
			\[
			g=
			\begin{pmatrix}
			a & b\\
			c & d
			\end{pmatrix}
			\in \SL_{2}(\mathbb{R})\text{ to }z\longmapsto g\cdot z:=\frac{az+b}{cz+d}
			\]
			is a group homomorphism $\SL_{2}(\mathbb{R})\rightarrow\operatorname{Bih}(\bbh^2)$, where $\operatorname{Bih}(\bbh^2)$ denotes the biholomorphic maps of the complex upper half plane $\bbh^2=\{z\in\mathbb{C}|\operatorname{Im}(z)>0\}$. Show that its kernel is 
			$\{\pm I \} $ where $I$ denotes as usual the $2 \times 2$ identity matrix.
			
			These maps are known as \emph{M\"{o}bius transformations}.
			
			%%%%%%% 1c solution
			
			\explanation{
			\begin{mdframed}
			\begin{solution*}
				Define the automorphy factor $j:\SL(2,\bbR) \times \mathbb{H}^2 \to \mathbb{C}$ by
				\begin{equation*}
					j(\gamma,z) = (cz+d) \quad \text{where } \gamma = \begin{pmatrix}a & b \\ c & d\end{pmatrix}
				\end{equation*}
				For all $z \in \mathbb{H}$ one has the trivial matrix relations
				\begin{equation*}
					\gamma \begin{pmatrix}z \\ 1\end{pmatrix} =  \begin{pmatrix}az+b \\ cz+d\end{pmatrix} = j(\gamma,z)  \begin{pmatrix}\gamma z \\ 1\end{pmatrix}
				\end{equation*}
				Given $\alpha,\beta \in \SL(2,\bbR)$ one now computes $\alpha \beta  \begin{pmatrix}z \\ 1\end{pmatrix}$ in two different ways: this yields both
				\begin{gather*}
				\alpha\beta\begin{pmatrix}z \\ 1\end{pmatrix} = j(\alpha\beta,z) \begin{pmatrix}(\alpha\beta)z \\ 1\end{pmatrix} \quad \& \quad \alpha\beta \begin{pmatrix}z \\ 1\end{pmatrix} = j(\alpha,\beta z)j(\beta,z)\begin{pmatrix}\alpha(\beta z) \\ 1\end{pmatrix}
				\end{gather*}
				
				It follows that we have the automorphy relation $j(\alpha\beta,z) = j(\alpha,\beta z)j(\beta,z)$ (also commonly referred to as the \emph{cocycle condition}) and furthermore that $(\alpha\beta)z = \alpha(\beta z)$. Hence $\SL(2,\bbR) \to \operatorname{Bih}(\mathbb{H}^2)$ is indeed a homomorphism, the remaining assertions are easy to verify.
			\end{solution*}
			\end{mdframed}
			}
			
			\item Prove that the induced homomorphism 
			\[
			\PSL_{2}(\mathbb{R})=\SL_{2}(\mathbb{R})/\left \{\pm I \right \}\rightarrow\operatorname{Bih}(\bbh^2)
			\]
			of (c) is actually an isomorphism. For the action of $\SL_{2}(\mathbb{R})$ on $\bbh^2$ from above determine the orbit $Gi$ and stabilizer $K$ of $i\in\bbh^2$. (Show also that $K$ is compact.) Using this, show that we have a diffeomorphism
			\[
			G/K\longrightarrow\bbh^2\text{, }g\longmapsto g\cdot i\text{.}
			\]
			
			%%%%%% 1d solution
			\explanation{
			\begin{mdframed}
			\begin{solution*}
				It suffices to show that any biholomorphism of $\mathbb{H}^2$ is actually a M\"{o}bius map. Note that the Cayley transform $\varphi:z \mapsto \tfrac{z-i}{z+i}$ maps the upper-half plane $\mathbb{H}^2$ biholomorphically to the (open) unit disc $\mathbb{D}$, and so we have an isomoprhism $\operatorname{Bih}(\mathbb{H}^2) \to \operatorname{Bih}(\mathbb{D}^2):\psi \mapsto \varphi \circ\psi\circ\varphi^{-1}$.
				
			As such, we identify the $\SL(2,\bbR)$ action on $\bbH^2$ with the $\operatorname{SU}_{1,1}(\mathbb{C})$ action on $\mathbb{D}$, and so it suffices to show that the latter is the group of biholomorphisms of $\mathbb{D}$	. The $\operatorname{SU}_{1,1}(\mathbb{C})$-action is given by
			\begin{equation*}
				\begin{pmatrix}a & b \\ \overline{b} & \overline{a}\end{pmatrix} \cdot z = \dfrac{az+b}{\overline{b}z+\overline{a}} \qquad \text{(where } |a|^2-|b|^2 = 1\text{)}
			\end{equation*}
			(You should verify this carefully!)
			
			
			Now let $\varphi$ be an arbitrary biholomorphism of $\mathbb{D}$, let $b = -\varphi(0)$ and set 
			\begin{equation*}
				\psi = \dfrac{1}{1+|b|^2}\begin{pmatrix} 1 & b \\ \overline{b} & 1\end{pmatrix} \circ \varphi
			\end{equation*}
			This is a biholomorphism of $\mathbb{D}$ which fixes $0$. The classical Schwarz lemma tells us that $\psi(z) = e^{i\theta}z$ for some $\theta \in [0,2\pi)$, which implies that in fact $\varphi \in \operatorname{SU}_{1,1}(\mathbb{C})$ whence the result. 
			
			To see that the action is transitive note that 
			\begin{equation*}
			\begin{pmatrix}y^{1/2} & xy^{-1/2} \\ 0 & y^{-1/2}\end{pmatrix} \cdot i = x+iy
			\end{equation*}
			for any $x+iy \in \mathbb{H}^2$. Furthermore,
			\begin{equation*}
				\begin{pmatrix}a & b \\ c & d\end{pmatrix}\cdot i = i  \Longleftrightarrow \dfrac{ai+b}{ci+d} = i \Longleftrightarrow ai+b = -c +id
			\end{equation*}
			Taking real and imaginary parts, one deduces that the stabiliser of $i$ is $\operatorname{SO}(2,\bbR)$. That $G/K \to \mathbb{H}^2$ is a diffeomorphism follows from standard arguments in Differential Geometry, see Helgason II.3.2 and II.4.3(a).
			\end{solution*}
			\end{mdframed}
			}
			
			\item Set $K=\SO_{2}(\mathbb{R})$,
			\[
			P= \left \{
			\begin{pmatrix}
			a & b\\
			0 & a^{-1}
			\end{pmatrix}
			~|~a,b\in\mathbb{R}\text{, }a>0 \right \},
			\]
			\[
			A= \left \{
			\begin{pmatrix}
			y^{1/2} & 0\\
			0 & y^{-1/2}%
			\end{pmatrix}
			~|~y\in\mathbb{R}^{+} \right \}\text{, and}%
			\]
			\[
			N = \left \{
			\begin{pmatrix}
			1 & x\\
			0 & 1
			\end{pmatrix}
			~|~x\in\mathbb{R} \right \}\text{.}
			\]
			
			Prove the Iwasawa decomposition, i.e.\ show that
			\[
			P\times K\longrightarrow G\text{, }(p,k)\longmapsto pk
			\]
			and
			\[
			N\times A\longrightarrow P\text{, }(n,a)\longmapsto na
			\]
			are diffeomorphisms. Are these also Lie group isomorphisms?\\ Show that $P$ is a semidirect product $N\rtimes A$ and that we have the diffeomorphism $N\times A\cong\bbh^2$.
			
			\emph{This decomposition is known as the \textbf{Iwasawa decomposition}.}
			%%%%%% 1e solution
			
			\explanation{
			\begin{mdframed}
			\begin{solution*}
			It is easy to verify that $P\cap K = N\cap A = \{\operatorname{Id}\}$, and so the maps $P\times K \to G$ and $N \times A \to P$ are injective.
			
			Consider now a matrix $g \in G$, we can consider it as determining a basis $g_i = g\cdot e_i$. Apply the Gram-Schmidt algorithm for the usual inner product on $\bbR^2$ to find a matrix $p \in P$ such that $p^{-1}g\in K$, a simple argument shows that $N\times A \to P$ is surjective. Neither of these maps are a homomorphism, however since $NA = P$, $N\cap A = \{1\}$ and $N \lhd P$ one has $P = N \rtimes A$.
			
			\emph{Caveat: Showing that a differentiable map is bijective does not suffice to prove that it is a diffeomorphism, however the Gram-Schmidt algorithm provides us directly with a differentiable inverse.}
			
			Finally $\bbh^2 \cong G/K \cong (N\times A\times K)/K \cong N\times A$. 
			\end{solution*}
			\end{mdframed}
			}
			
			\item Prove that $K$ is unimodular by showing that $d\mu_{\bbh^2}=y^{-2}dxdy$, $z=x+iy$, is a $G$-invariant volume form on the $G$-homogeneous space $\bbh^2$.
			
			%%%%%% 1f solution
			
			\explanation{
			\begin{mdframed}
			\begin{solution*}
				The derivative of $f_g:z \mapsto \tfrac{az+b}{cz+d}$ is
				\begin{equation*}
					\dfrac{a(cz+d)-c(az+d)}{(cz+d)^2} = (cz+d)^{-2}
				\end{equation*}
				Write $\operatorname{Re}(f_g) = u_g$ and $\operatorname{Im}(f_g) = v_g$. The Cauchy-Riemann equations imply that 
				\begin{equation*}
					\begin{pmatrix} \tfrac{\partial u_g}{\partial x} &  \tfrac{\partial u_g}{\partial y} \\  \tfrac{\partial v_g}{\partial x} &  \tfrac{\partial y_g}{\partial x}
					\end{pmatrix} = \begin{pmatrix} \operatorname{Re}(cz+d)^{-2} & -\operatorname{Im}(cz+d)^{-2} \\ \operatorname{Im}(cz+d)^{-2} & \operatorname{Re}(cz+d)^{-2} \end{pmatrix}
				\end{equation*}
				and the determinant $\Delta$ of this matrix is $|cz+d|^{-4}$. Now we calculate
				\begin{equation*}
					g^*(dxdy) = \Delta dxdy = |cz+d|^{-4} \quad \text{and}\quad g^*(y^{-2}) = \operatorname{Im}\left(\dfrac{az+b}{cz+d}\right)^{-2} = y^{-2}|cz+d|^{-4}
				\end{equation*}
			and therefore $g^\mu_{\bbh^2} = \mu_{\bbh^2}$ and the assertion follows from Weil's formula.
			\end{solution*}
			\end{mdframed}
			}
			
			\item Let \[\mathcal{F}\coloneqq\{z\in\bbh^2~|~(|z|>1\text{ and } -1/2\leq\operatorname{Re}(z)<1/2) \text{ or }(|z|=1\text{ and }-1/2\leq\operatorname{Re}(z)\leq 0)\}.\]
			Show that for all $z\in \bbh^2$ the orbit $G_\mathbb{Z}z$ intersects $\mathcal{F}$ in a unique point. 
			
			\hint \ For every $G_{\mathbb{Z}}$-orbit $G_{\mathbb{Z}}z$, $z\in\bbh^2$, consider $w\in G_{\mathbb{Z}}z$ with maximal imaginary part.
			
			%%%%%% 1 g solution
			
			\explanation{
			\begin{mdframed}
			\begin{solution*}
			Note that 
			\begin{equation*}
				\operatorname{Im}\left(\dfrac{az+b}{cz+d}\right) = \dfrac{\operatorname{Im}(z)}{|cz+d|^2}  \begin{cases} = \operatorname{Im}(z)|d|^{-2} & \text{if } c=0 \\ \leq  \operatorname{Im}(z)|c|^{-2} & \text{otherwise} \end{cases}
			\end{equation*}
			Hence given some $z \in \bbh^2$ the function $G_{\mathbb{Z}} \to \bbh:\gamma \mapsto \operatorname{Im}(\gamma z)$ obtains a maximum, say at $w = \gamma_0 z$. Since $(\begin{smallmatrix} 1 & 1 \\ 0& 1\end{smallmatrix})$ acts by translations $z \mapsto z+1$ we may assume that $-1/2 \leq \operatorname{Re}(w) < 1/2$. In addition since $(\begin{smallmatrix} 0 & 1 \\ -1 & 0\end{smallmatrix})$ acts as inversion $z \mapsto z^{-1}$ and $\operatorname{Im}(-1/w) = \operatorname{Im}(w)|w|^{-2}$ one clearly has $|w| \geq 1$.
			
			It remains to show that we may impose $-1/2 \leq \operatorname{Re}(w) \leq 0$ if $|w| = 1$ and this follows from considering $z \mapsto -1/z$ again.
			
			Thus each orbit $G_{\mathbb{Z}}\cdot z$ intersects $\mathcal{F}$, to show it is a fundamental domain it remains to show that $z,\gamma z \in \mathcal{F}$ implies $z = \gamma z$ --- this follows by similar considerations to the above, and is left to the reader.
			\end{solution*}
			\end{mdframed}
						
			}
			
			\item Show that the volume of $\mathcal{F}$ with respect to $\mu_{\bbh^2}$ is $\pi/3$. Deduce that $\mu(G/G_{\mathbb{Z}})<\infty$.
			
			%%%%%%% 1h solution
			
			\explanation{
			\begin{mdframed}
			\begin{solution*}
			We can calculate
			\begin{align*}
			\int_{\mathcal{F}} \, d\mu_{\bbh^2} & = \int_{-1/2}^{1/2} \int_{\sqrt{1-x^2}}^\infty y^{-2}\, dx \\
			& = \int_{-1/2}^{1/2}  -y^{-1}\big|_{y=\sqrt{1-x^2}}^{y=\infty}\, dydx \\
			& = \int_{-1/2}^{1/2}(1-x^2)^{-1/2}\, dx \\
			& = \arcsin(x)\big|_{x=-1/2}^{x=1/2} = \pi/3.
			\end{align*}
			The second assertion follows from part (b) applied to $A = \mathcal{F}$.
			\end{solution*}
			\end{mdframed}

			
			}

		\end{enumerate}
		
%%%%%%%%%%%%%%%%%%%%% QUESTION 2 %%%%%%%%%%%%%%%%%%%%%%%%%%

\item[2.] Consider the hyperbolic $n$-space
$$
\bbh^n = \left\{ p \in \bbr^{n+1} \colon b(p,p)=-1 \text{ and } p_{n+1} > 0 \right\}
$$
defined by the bilinear form $b(p,q)=p_1q_1 + \ldots + p_n q_n - p_{n+1}q_{n+1}$. The tangent space at a point $p\in \bbh^n$ is defined as
$$
T_p\bbh^n = \left\{ x \in \bbr^{n+1} \colon 
\begin{matrix}
\text{ There exists a smooth path } \gamma \colon (-1,1) \to \bbh^n \\
 \text{ such that } \gamma(0)=p \text{ and } \dot{\gamma}(0)=x
\end{matrix} \right\}.
$$

\begin{enumerate}
	%%%%%% 2a
	\item Show that $T_p\bbh^n = \left\{ x \in \bbr^{n+1} \colon b(p,x)=0 \right\}$.
	
	%%%%%% 2a solution
	\explanation{\begin{mdframed}\begin{solution*}
		Consider any $x\in T_p\mathbb{H}^n$, and let $\gamma \colon (-1,1)\to \mathbb{H}^n$ be a smooth path such that $\gamma(0)=p$ and $\dot{\gamma}(0)=x$. For every $t \in (-1,1)$, $b(\gamma(t),\gamma(t))=-1$, since $\gamma$ takes values in $\mathbb{H}^n$. We write $\gamma(t)=(\gamma_1(t),\cdots,\gamma_{n+1}(t))$. 
		
		Taking the derivative we see that
    $$
    0=\frac{d}{dt}b(\gamma(t),\gamma(t))= \frac{d}{dt}\left(\sum_{i=1}^n \gamma_i(t)^2 - \gamma_{n+1}(t)^2\right) = \sum_{i=1}^n 2 \gamma_i(t) \dot{\gamma}_i(t) - 2 \gamma_{n+1}(t) \dot{\gamma}_{n+1}(t) 
    $$
    and at $t=0$ this is
    $$
    0=\sum_{i=1}^n \gamma_i(0)  \dot{\gamma}_i(0) - \gamma_{n+1}(0) \dot{\gamma}_{n+1}(0) = \sum_{i=1}^n p_i x_i - p_{n+1}\cdot x_{n+1}  = b(p,x).
    $$
	Thus we have shown that $T_p\mathbb{H}^n \subset \left\{ x \in \bbR^{n+1} \colon b(p,x)=0 \right\}$, but since $\dim{T_p\mathbb{H}^n}=n$ we must have equality.
	\end{solution*}
	\end{mdframed}}
	
	%%%%%% 2b
	\item Show that $g_p = b|_{T_p\bbh^n} \colon T_p\bbh^n \times \, T_p\bbh^n \to \bbr$ is a positive definite symmetric bilinear form on $T_p\bbh^n$ (this means that $g_p$ is a scalar product, and $(\bbh^n, g)$ is a Riemannian manifold).
	
	\emph{Hint: Use (a) and the Cauchy-Schwarz-inequality on $\bbr^n$.}
	
	%%%%%% 2b solution
	\explanation{\begin{mdframed}\begin{solution*}Bilinearity and symmetry $b(x,y)=b(y,x)$ follow directly, so we just need to show positive definiteness. Using (a) we can 		write any $p\in \bbh^n \subset \bbr^n \times \bbr$ and $x\in T_p\bbh^n \subset \bbr^n\times \bbr$ as
		$$
	p=\left( p', \sqrt{|p'|^2+1}\right) \in \bbH^n\subset \bbR^n\times \bbR 
	$$
	$$
	x=\left( x', \frac{\langle p',x'\rangle}{\sqrt{|p'|^2+1}}\right) \in T_p\bbH^n \subset \bbR^n \times \bbR
	$$ 
	where $\langle\cdot \ , \cdot \rangle$ is the standard scalar product in $\bbR^n$. To show positive definiteness it remains to prove that for all $x \in T_p\bbH^n$
	$$
	b(x,x)  \geq 0.
	$$
	with equality if and only if $x = 0$.
	Indeed, by the Cauchy-Schwarz-inequality 
	$$\langle p',x'\rangle^2 \ \leq |p'|^2  |x'|^2\leq |p'|^2|x'|^2+|x'|^2 = (|p'|^2+1)|x'|^2$$
	and thus
	$$
	|x'|^2 \geq \frac{\langle p',x'\rangle^2}{|p'|^2+1}
	$$
	and so
	$$
	b(x,x)= |x'|^2 - \frac{\langle p',x'\rangle^2}{|p'|^2+1} \geq 0
	$$
	with equality if and only if $x' = 0$ (which forces $x=0$).
	\end{solution*}
	\end{mdframed}}
	
	%%%%%% 2c
	\item Show that the map $s_p \colon \bbr^{n+1}\to \bbr^{n+1}, q \mapsto -2p \cdot b(p,q)-q$ defines a well defined geodesic symmetry of $\bbH^n$ (that is, it is an  involution with an isolated fixed point $p$). This means that the hyperbolic plane $\bbh^n$ is a Riemannian (globally) symmetric space.
\end{enumerate}	

%%%%%% 2c solution
\explanation{\begin{mdframed}\begin{solution*}
	\newcommand{\wur}[1]{ \sqrt{|\overrightarrow{#1}|^2+1} }
		%%%%% s_p well-defined
	To see firstly that $s_p$ is well-defined we write as before
	$$
	p=\left(  p', \sqrt{|p'|^2+1}\right), \quad q=\left( q',\sqrt{|q'|^2+1}\right) \in \bbH^n\subset \bbR^n\times \bbR .
	$$
	and so in particular
	\begin{equation*}
		b(p,q) = \langle p',q'\rangle - \sqrt{|p'|^2+1}\sqrt{|q'|^2+1}
	\end{equation*}
	By using the definition we have that
	\begin{align*}
		s_p(q) & = -2p\cdot b(p,q) - q \\
		& =  \left(-2p'\cdot b(p,q) - q', -2\sqrt{|p'|^2+1}\cdot b(p,q) - \sqrt{|q'|^2+1}\right)
	\end{align*}
	and so we calculate
	\begin{align*}
		b(s_p(q),s_p(q)) & = 4|p'|^2b(p,q)^2 + 4\langle p',q'\rangle b(p,q) + |q'|^2 \\ &- \Big(4(|p'|^2+1)b(p,q)^2 + 4\sqrt{|p'|^2+1}\sqrt{|q'|^2+1} \cdot b(p,q) + |q'|^2-1\Big) \\
		& = 4\langle p',q'\rangle b(p,q) - 4b(p,q)^2 - 4\sqrt{|p'|^2+1}\sqrt{|q'|^2+1}\cdot b(p,q) -1\\
		& = 4b(p,q)b(p,q)-4b(p,q)^2-1 = -1
	\end{align*}
	So indeed $s_p(q) \in \bbh^n$ as required.
		
	Note also that $s_p(p)=-2p(-1)-p=p$ is a fixed point.
	%%%%% s_p is an isometry
	Next we show that $s_p$ is an isometry, so we need to look at the differential
	\begin{align*}
	d_p s_p \colon T_pM &\to T_{s_p(p)}M=T_p M .
	\end{align*}
	If we write the points $q,p \in \bbH^n \subset \bbR^{n+1}$ in the standard basis $\{e_i\}_i$, we get the partial derivatives
	\begin{align*}
	\frac{\partial}{\partial x_i} b(p,\cdot ) &=\left\{ \begin{matrix} p_i \quad \text{if } i \leq n \\ -p_{n+1} \quad \text{if } i=n+1 \end{matrix}\right. \\
	\frac{\partial}{\partial x_i} s_p &= \left\{ \begin{matrix} -2p\cdot p_i - e_i \quad \text{if } i \leq n \\ 2p\cdot p_{n+1} - e_{n+1} \quad \text{if } i=n+1 \end{matrix}\right.
	\end{align*}
	and thus for $v \in T_p M$ we have
	\begin{align*}
	(d_ps_p) v &= \left( \begin{matrix}
	-2p_1^2-1 && -2p_1p_2 && \cdots && 2p_1p_{n+1} \\
	-2p_2p_1 && -2p_2^2 -1 && \cdots && 2p_2p_{n+1} \\
	\vdots && \vdots && \ddots && \vdots \\
	-2p_{n+1}p_1 && -2p_{n+1}p_2 && \cdots && 2p_{n+1}^2-1
	\end{matrix} \right) v \\
	&=  \left( \begin{matrix}
	-2p_1^2v_1 -2p_1p_2v_2 - \cdots + 2p_1p_{n+1}v_{n+1} \\
	-2p_2p_1v_1-2p_2^2v_2 - \cdots + 2p_2p_{n+1}v_{n+1} \\
	\vdots  \\
	-2p_{n+1}p_1v_1 -2p_{n+1}p_2v_2 - \cdots + 2p_{n+1}^2v_{n+1}
	\end{matrix} \right)  - v \\
	&=  \left( \begin{matrix}
	-2b(p,v)p_1 \\
	-2b(p,v)p_2 \\
	\vdots  \\
	-2b(p,v)p_n
	\end{matrix} \right)  - v = - v
	\end{align*}
	where we used that $b(p,v)=0$ from part (a). Using linearity we have that
	$$
	g_{s_p(p)}((d_ps_p)v,(d_ps_p)w)=g_p(-v,-w)=g_p(v,w),
	$$
	so $s_p$ is an isometry.
	%%%%% Isolated fixed point
	
	Now let us argue why $p$ is an isolated fixed point, indeed suppose that $s_p(q) = q$ for some $q\in \bbh^n$. Then $-2p\cdot b(p,q) - q = q$, so $q = -b(p,q)p$, and in particular $q = \lambda p$ for $\lambda$ some constant. Then $-1 = b(q,q) = b(\lambda p, \lambda p) = \lambda^2b(p,p) = -\lambda^2$ and so $\lambda = \pm 1$. $\lambda = -1$ corresponds to $q_{n+1} <0$, which is excluded since $\bbh^n$ is only the upper sheet of the hyperboloid, and so $q=p$ and $s_p$ has an isolated fixed point.
	
	%%%%% Geodesic symmetry	
	 By lemma II.5 of the lecture, $d_ps_p = -\operatorname{Id}_{T_p\bbH^n}$ is equivalent to $s_p \circ s_p = \operatorname{Id}_{\bbH^n}$, and so we are done. Alternatively we can calculate
	\begin{align*}
	    s_p\circ s_p (q)&= s_p (-2p\cdot b(p,q)-q) \\
	    &= -2p\cdot b(p,-2p\cdot b(p,q)-q)-(-2p\cdot b(p,q)-q) \\
	    &= 4p\cdot b(p,q)b(p,p)+2p\cdot b(p,q)+2p\cdot b(p,q)+q = q
	\end{align*}
\end{solution*}	\end{mdframed}}

%%%%%%%%%%%%%%%%%% QUESTION 3 %%%%%%%%%%%%%%%%%%%

\item[3.]Show that $A\mapsto gAg^t$ defines a group action of $\operatorname{SL}(n,\bbr) \ni g$ on
$$
\mathcal{P}^1(n)=\left\{A \in M_{n\times n}(\bbr) \colon A=A^t ,\  \det A=1 ,\ \ A>\!>0 \right\}.
$$
Show that this action is transitive (that is, $\forall A,B \in \mathcal{P}^1(n)$ there exists some $ g \in \operatorname{SL}(n,\bbr)$ such that $ gAg^t=B$). 

\emph{You may use the Linear Algebra fact that symmetric matrices are orthogonally diagonalisable (that is, if $A=A^t$, then $\exists Q \in \operatorname{SO}(n,\bbr)$ such that $QAQ^t$ is diagonal)}.

\explanation{\begin{mdframed}\begin{solution*}
We write the group action as $g.A=gAg^t$. We first need to show that the action is well defined.

\begin{itemize}
	\item \emph{(Symmetry)} $(g\cdot A)^T = (gAg^T)^T = gA^Tg^T = gAg^T = g\cdot A$;
	\item \emph{(Determinant)} $\det(g\cdot A) = \det(g)\det(A)\det(g^T) = \det(A) = 1$;
	\item \emph{(Positive definite)} Let $x\in \bbR^n\backslash \{0\}$, then $$x^T(g\cdot A)x = x^TgAg^Tx = (g^Tx)^T A (g^Tx) > 0$$ since $g^Tx \in \bbR^n\backslash\{0\}$.
\end{itemize}

\noindent
We also note that it is a group action --- clearly $\operatorname{Id}\cdot A = A$, and $(gh)\cdot A = ghAh^Tg^T = g\cdot(h\cdot A)$.

\noindent
It remains to show that the action is transitive. Let $A,B\in \mathcal{P}^1(n)$. We can use linear algebra to get $Q,R \in \operatorname{SO}(n)<\operatorname{SL}(n,\bbR)$ such that $Q.A$ and $R.B$ are diagonal, have determinant 1 and are positive definite (by the well-definedness of the group action). Positive definiteness implies that all entries are non-negative. Then the matrix $\Lambda=(Q\cdot A) (R\cdot B)^{-1}$ is also diagonal, has determinant 1 and positive elements on the diagonal. We can therefore take the component-wise square root $\sqrt{\Lambda}$ of $\Lambda$. 

Set $g=Q^{-1}\sqrt{\Lambda}R\in\operatorname{SL}(n,\bbR)$ and use the fact that $R.B$ commutes with $\sqrt{\Lambda}$ since they are diagonal to see that
\begin{align*}
g\cdot B&=Q^{-1}\sqrt{\Lambda}R\cdot B=Q^{-1}\cdot \left(\sqrt{\Lambda}(R\cdot B){\sqrt{\Lambda}}^T\right) =Q^{-1}.\left(\sqrt{\Lambda}{\sqrt{\Lambda}}^T  R\cdot B\right) \\ &= Q^{-1}\cdot(\Lambda R\cdot B)
= Q^{-1}\cdot ((Q\cdot A)(R\cdot B)^{-1}(R\cdot B))=Q^{-1}\cdot Q \cdot A=A.
\end{align*}
this shows that from any point $B\in \mathcal{P}^1(n)$ you can go to any point $A \in \mathcal{P}^1(n)$ by the action of $\operatorname{SL}(n,\bbR)$ (the action is transitive). 
\end{solution*}
\end{mdframed}}
\end{enumerate}

\end{document}